Quand on se pose la question de savoir ce qu'est une bibliothèque d'orchestre, plusieurs éléments peuvent être examinés, tels que la différence avec les bibliothèques musicales ou encore le métier spécifique de musicothécaire.

\subsection*{Le rôle traditionnel des bibliothèques musicales et d'orchestre} %le package section* permet d'annoncer une partie de chapitre sans la numéroter.
Les bibliothèques d'orchestre sont un type particulier de bibliothèques musicales. Comme toute bibliothèque, ce sont des \enquote{lieux de conservation et de transmission de la connaissance}%enquote est un package pour inclure des guillemets 
\footnote{Carbone (Pierre), Introduction. Dans : \textit{Les bibliothèques}. %package pour mettre un contenu en italique 
Paris: Presses Universitaires de France. Que sais-je ?, 2017.p.3-10.
% package pour créer une note de bas de page. 
URL:\url{https://shs-cairn-info.portail.psl.eu/les-bibliotheques--9782130787549-page-3?lang=fr}}. %package pour inclure une url. 
Cependant, elles diffèrent, au niveau de leurs missions, des bibliothèques traditionnelles que sont les bibliothèques municipales par exemple.
En effet, les bibliothèques d'orchestre sont, comme leur nom l'indique, avant tout au service d'un orchestre. Leur raison d'être réside dans la préparation des productions de la saison musicale. 

Excepté l'ouvrage de Gilles Pierret\footnote{Pierret (Gilles), dir. \textit{Musique en bibliothèque}, Paris, Éditions du Cercle de la librairie, 2012, 357 p. (« Bibliothèques ». Ouvrage de 2002 réédité en 2012).}, \textit{Musique en bibliothèque}, la littérature scientifique francophone a peu traité le sujet des bibliothèques d'orchestre et de leurs bibliothécaires.
Si le concept de bibliothèque musicale a été défini par l'ouvrage de Gilles Pierret, ancien directeur de la Médiathèque musicale de Paris, il n'existe pas d'ouvrage de référence sur la bibliothèque d'orchestre qui pourrait en présenter les caractéristiques et les missions.
C'est dans la littérature anglophone qu'il faut chercher ce type d'ouvrage.
Par exemple, Russ Girsberger proposait dans l'introduction du manuel de référence "\textit{A Manual for the Performance Library}" une définition de la bibliothèque d'orchestre :
\begin{quote}
	The performance library is the repository and distribution center for music scores and parts used by performing musicians and ensembles. Such libraries support many types of musical groups, including symphony orchestras, opera companies, ballet companies, wind ensembles and bands, choruses, chamber ensembles, and jazz and popular music combos.\footnote{« Une bibliothèque d'orchestre est un lieu de conservation et de diffusion des partitions et des parties séparées utilisées par les musiciens et l\textit{}es ensembles. Ces bibliothèques accompagnent de nombreux types de formations musicales, notamment les orchestres symphoniques, les compagnies d'opéra, les compagnies de ballet, les ensembles à vent et les fanfares, les chœurs, les ensembles de musique de chambre, ainsi que les groupes de jazz et de musique populaire. », Girsberger (Russ), A Manual for the Performance Library, Lanham, Scarecrow Press, 2006, 174 p.}
\end{quote} 

Pour définir la bibliothèque d'orchestre, il paraît intéressant d'étudier le rôle des bibliothécaires d'orchestre. Cette profession méconnue, mais essentielle à la vie de l'orchestre, est aussi dénommée partothécaire, musicothécaire ou encore\textit{ performance librarian} dans le monde anglo-saxon.
Dans sa revue critique de l'ouvrage "\textit{Stories and Lessons from the World's Leading Opera, Orchestra Librarians, and Music Archivists, Volume 2 : Europe and Asia}"\footnote{"Lo (Patrick), Sutherland (Robert), Hsu (Wei-En) et Girsberger (Russ), Stories and Lessons from the World's Leading Opera, Orchestra Librarians, and Music Archivists, Volume 1:North and South America \& Volume 2 : Europe and Asia, Bingley, Emerald Publishing Limited, 2022, 780 p.}, Geneviève Beaudry de la bibliothèque Marwin Duchow de l'Université de McGill décrit ainsi le métier de bibliothécaire d'orchestre :
\begin{quote}
	Les musicothécaires (en anglais\textit{ Performance Librarians}) forment un groupe un peu à part. Ces individus passionnés sont employés par des orchestres, maisons d'opéra et de ballet, harmonies militaires ou institutions d'enseignement. Ces personnes travaillent souvent seules. Ils ont la plupart du temps appris leur métier avec des collègues plus expérimentés qu'eux. Aucune formation professionnelle n'existe pour ce métier. Leur rôle consiste principalement à acquérir et préparer les partitions d'ensemble pour tous, afin que chef, instrumentistes, solistes et choristes aient en main le matériel nécessaire à leurs prestations. Ce sont aussi eux qui s'occupent du catalogage des collections, et de la gestion des droits de performance pour les œuvres encore sous copyright.
\end{quote}

Cette définition éclaire nombre d'aspects de la profession de bibliothécaire d'orchestre. Professionnel spécialisé, en charge à la fois de la préparation des concerts et de fonctions bibliothéconomiques telles que le catalogage, le bibliothécaire d'orchestre est autant gestionnaire que musicien, bibliothécaire que musicologue. Il est d'ailleurs parfois appelé bibliothécaire-musicologue, comme à l'orchestre \textit{Insula Orchestra}, orchestre de pratique historiquement informée en résidence à la Seine Musicale à Boulogne.
Véritabe couteau suisse, le bibliothécaire d'orchestre se doit d'être soucieux des détails et d'être capable de gérer plusieurs tâches de front tout en respectant des délais parfois serrés.\footnote{Karen Schnackenberg, Principal Librarian, Dallas Symphony Orchestra dans un entretien restitué dans l'ouvrage Lo (Patrick), Sutherland (Robert), Hsu (Wei-En) et Girsberger (Russ), Stories and Lessons from the World's Leading Opera, Orchestra Librarians, and Music Archivists, Volume 1:North and South America \& Volume 2 : Europe and Asia, Bingley, Emerald Publishing Limited, 2022, 780 p.}.

Pourtant, le bibliothécaire a souvent appris son métier de façon empirique, en raison d'une faible offre de formation. 
En France, il existe une formation reconnue de bibliothécaire d'orchestre, au CFA Métiers des Arts de la Scène de Lorraine. Il s'agit de la licence professionnelle \enquote{Communication et valorisation de la création artistique} parcours \enquote{Métiers de la scène lyrique} proposée en partenariat avec l'Université de Lorraine.  Les étudiants y choisissent l'option \enquote{bibliothécaire musical} et suivent une année de scolarité en alternance dans un orchestre ou opéra partenaire.\footnote{Pour découvrir la formation du CFA Métiers des Arts de la Scène, voir la vidéo du CFA des Arts de la Scène de Lorraine. Maud Rolland, Apprentissage en tant que bibliothécaire d'orchestre à Radio France, vidéo YouTube, CFA Métiers des arts de la scène, 2019, disponible sur \url{https://www.youtube.com/watch?v=oseackOdb2A} (consulté le 8 juillet 2026).}

Une lecture de fiches de poste des recrutements de bibliothécaires d'orchestre (à Toulouse\footnote{Voir annexe \ref{annexe1toulouse.pdf} page 42}%pour faire un renvoi vers l'annexe, 
ou encore Monte-Carlo\footnote{Voir annexe \ref{annexe2montecarlo.pdf} page 44}) nous apprend que les missions traditionnelles du bibliothécaire d'orchestre sont très variées mais toutes tournées vers l'objectif des productions musicales.

En effet, le bibliothécaire s'assure de numériser tout le « matériel », c'est-à-dire l'ensemble des partitions des musiciens, en cas de perte, mais également avant et après répétition, afin de garder une traçabilité de l'exécution des œuvres. Cela représente un intérêt tout particulier pour les ensembles adeptes de la pratique historiquement informée (HIP) où les chefs et musiciens sont sensibles à la tradition historique de l'interprétation d'une œuvre et soucieux d'en documenter le processus de production.

\subsection*{De nouvelles perspectives pour les bibliothèques d'orchestre :  un réseau au service de la science ouverte}

Depuis les années 2000, les bibliothèques sont en profonde mutation en raison du numérique. 
Pierre Carbone le présentait ainsi dans la conclusion de son ouvrage Les bibliothèques (collection « Que sais-je ?»)\footnote{Carbone (Pierre), \textit{Introduction}. Dans \textit{Les bibliothèques}. Paris: Presses Universitaires de France. Que sais-je ?, 2017, p.3-10. URL:\url{https://shs-cairn-info.portail.psl.eu/les-bibliotheques--9782130787549-page-3?lang=fr}}:

\begin{quote}
	Lieux de mémoire physique mais aussi virtuelle, les bibliothèques sont ici et maintenant des ressources vives pour leurs usagers, que ceux-ci viennent sur place ou y accèdent à distance. [...] Les usages et les services se diversifient, de nouveaux modèles apparaissent, les modes d'organisation et de fonctionnement des bibliothèques évoluent. [Les bibliothèques sont] à la fois lieux physiques et services numériques présents sur les réseaux.
\end{quote}

Le constat de Pierre Carbone est également valable pour les bibliothèques musicales et d'orchestre que le numérique a profondément transformées. Comme les bibliothèques classiques, de nombreuses bibliothèques musicales proposent en effet un service numérique comme La Philharmonie à la demande qui est la bibliothèque numérique de la Philharmonie de Paris.

 Ainsi, si la mission des bibliothèques classiques, comme musicales, reste la conservation et la valorisation de leurs collections sur le long terme, l’émergence du numérique les transforme en profondeur et diversifie leurs missions et services. Aujourd’hui, les bibliothèques sont des lieux physiques à part entière mais proposent également, dans la majorité des cas, une offre de services en ligne. Au gré de leur présence sur les réseaux sociaux ou sur les réseaux professionnels de recherche ou des bibliothèques, leur identité évolue même si l’objectif premier reste de « donner librement accès à la connaissance, à l’information et à la culture. » (Pierre Carbone, 2002)
Parallèlement à cette mutation numérique, la dimension sociale et pédagogique des bibliothèques se renforce et, dans un contexte d’essor de l’intelligence artificielle et de l’usage croissant des outils numériques, les missions traditionnelles des bibliothèques de conservation et de valorisation patrimoniale se doublent d’une nécessité de  « développer les compétences informationnelles au sein de la société»\footnote{CARBONE (Pierre), \textit{Conclusion}, dans \textit{Les bibliothèques}. Paris: Presses Universitaires de France. Que sais-je ?, 2017, p.123.}.

Ainsi, les bibliothèques s’affirment comme piliers de la démocratie du savoir : «La lutte contre les manipulations de l’information, la souveraineté et la fiabilité des informations et des données, la science ouverte et la création d’écosystème numérique en sources ouvertes, ou encore la maîtrise de l’usage des IA génératives» sont autant d’enjeux sociétaux qui relèvent directement des compétences des bibliothèques. (Laure Darcos)\footnote{Darcos (Laure), Mieux reconnaître le rôle des conservateurs et conservateurs généraux des bibliothèques, Question écrite n° 08054, 17ᵉ législature, publiée au JO Sénat le 19 mars 2026, page 1354.} 

La coopération entre bibliothèques, favorisée par le numérique, leur permet de s’ancrer dans ces questions sociétales, notamment en renforçant le lien avec le système universitaire, dans un contexte d’essor de la science ouverte. Historiquement, cette coopération entre bibliothèques a débuté par le partage des techniques des bibliothèques (mutualisation bibliographique) et l’adoption, par la majorité des bibliothèques publiques, de la classification décimale Dewey (fin du XIX$^{e}$ siècle) et son équivalent en BU : la classification décimale universelle (XX$^{e}$ siècle). 

L’informatisation des catalogues a permis de faire des recherches à critères multiples et des recherches dans les bases de données bibliographiques. 

A l’enjeu de normaliser la classification (le classement des fonds) au XIX$^{e}$ et XX$^{e}$ siècles a succédé l’enjeu de l’indexation et de la recherche plein texte dans les catalogues.  
Au départ, il s’agissait de localiser physiquement un livre dans un réseau de bibliothèques. Aujourd’hui, il s’agit de rendre les données documentaires interopérables. 
Cette interopérabilité des données documentaires est tout autant recherchée par les bibliothèques musicales et d'orchestre pour la conservation de leur patrimoine. Comme les bibliothèques classiques, les bibliothèques d'orchestre veulent transformer leur patrimoine en données exploitables à la fois pour la recherche et pour l'interprétation musicale, notamment l'interprétation musicalement informée. 

Ces enjeux d'interopérabilité sont héritiers des avancées et des politiques du numérique en bibliothèques qui sont à l'œuvre depuis les années 1960. 
En effet, à  la fin des années 1960, la structuration internationale bibliographique émerge. Le format MARC (\textit{machine readable cataloguing}) initie le catalogage informatisé. L’IFLA (Fédération internationale des associations et institutions de bibliothécaires) encourage la diffusion de ce modèle (UNIMARC, MARC21) afin de permettre l’échange des notices à l’échelle mondiale. 

En France, ces enjeux de réseau bibliographique s’incarnent avec la création de deux catalogues : le catalogue du Sudoc  (Système universitaire de documentation) géré par
l’ABES (Agence bibliographique de l’enseignement supérieur) et le CCFr (Catalogue collectif de France) géré par la BnF.

Les bibliothèques musicales aussi ont créé un équivalent au Sudoc, le catalogue du réseau des bibliothèques de la CMF\footnote{Le Catalogue du réseau des bibliothèques de la CMF est disponible à cette adresse : \url{https://cmf-documentation.org/}} (Confédération Musicale de France)

Cependant, le Web souligne les limites des formats bibliographiques traditionnels. Ceux-ci avaient été conçus pour informatiser les fiches cartonnées des catalogues et pour les cantonner dans des bases de données isolées. 
Pour y remédier, l’IFLA a développé le modèle IFLA-LRM (\textit{Library Reference Model}), publié en 2018 et concrétisé par le code de catalogage RDA (Ressources : Description et Accès) pour structurer les données des bibliothèques sur le Web. Avec ce modèle, l’IFLA permet de dépasser l’enjeu d’échange de notices pour permettre d’inclure les données documentaires dans le Web de données (\textit{Linked Open Data}).

L’objectif n’est plus de gérer un catalogue fermé mais de transformer chaque entité (auteur, oeuvre, édition) en donnée ouverte, interopérable et identifiée par des référentiels d’autorité comme IdRef (Identifiants et Référentiels pour l’enseignement supérieur et la recherche), ISNI (\textit{International Standard Name Identifier}), ou Wikidata. Ainsi, les bibliothèques peuvent lier leurs fonds aux données de la recherche. Ce basculement vers le Web de données constitue le socle de la science ouverte puisqu’il permet aux chercheurs d’accéder aux données de la recherche. 
C’est précisément pour cette raison qu’a été créé le consortium d’orchestres LaDocumenta.eu\footnote{Pour une explication plus détaillée de LaDocumenta.eu, voir l'\nameref{subsec:ladocumenta_cas}}, s'inspirant d'Europeana, mais centrée sur les bibliothèques d'orchestre de pratique historiquement informée, cette plateforme permet un partage de données de la recherche dans une logique de science ouverte.

\begin{quotation}La science ouverte est la diffusion sans entrave des résultats, des méthodes et des produits de la recherche scientifique. Elle s’appuie sur l’opportunité que représente la mutation numérique pour développer l’accès ouvert aux publications et – autant que possible – aux données, aux codes sources et aux méthodes de la recherche\footnote{(Ministère de l’Enseignement supérieur et de la recherche\url{https://www.enseignementsup-recherche.gouv.fr/fr/une-politique-ambitieuse-pour-la-science-ouverte-50735})}.\end{quotation} 

Ce concept s’appuie sur des textes fondateurs datant des années 2000 : les déclarations internationales de Budapest en 2002 et de Berlin en 2003 ont souligné les besoins théoriques d’un accès libre à la recherche, en parallèle de l'ouverture de la plateforme HAL, une plateforme pluridisciplinaire de dépôt de textes académiques d’archives ouvertes. En outre, la Loi pour une République numérique de 2016 donne un cadre légal français de la science ouverte en France en consacrant le droit pour les chercheurs de diffuser en accès libre leurs travaux financés par des fonds publics.

Les bibliothèques d’orchestre sont également concernées par les transformations des métiers des bibliothèques en raison de l’essor du numérique.
Historiquement, les bibliothèques d’orchestre étaient centrées sur la conservation, la gestion matérielle des fonds ainsi que sur le suivi des productions de concerts. Avec l’émergence du numérique, ces missions demeurent mais la philosophie du métier a évolué. En effet, les bibliothécaires d’orchestre s’ancrent à présent dans une logique de travail en réseau et de coopération entre bibliothèques, parfaitement similaire aux changements opérés dans les bibliothèques classiques que nous venons de décrire. 

Si les catalogues circulent plus aisément, le suivi documentaire des productions est également facilité grâce aux outils numériques. 
Dans le domaine de la pratique historiquement informée (HIP, en anglais \textit{Historically Informed Performance}) en particulier, ces technologies numériques transforment le travail artistique et scientifique des musiciens et des bibliothécaires. Le travail artistique, centré sur l’exécution musicale bénéficie pleinement de ces outils pour s’enrichir d’un travail scientifique assisté du numérique. 

En effet, l’accès aux sources, fac-similés et éditions critiques conservés en bibliothèque d’orchestre est grandement facilité depuis que leur catalogue est public. Parallèlement, ces ressources des bibliothèques deviennent un objet d’étude phare des chercheurs intéressés par la pratique historiquement informée. C’est ainsi que, mus par une même volonté d’accéder aux sources musicales inédites des bibliothèques, les chercheurs et les musiciens ayant une pratique historiquement informée constituent ensemble des réseaux professionnels afin de travailler dans une logique de science ouverte. La bibliothèque devient alors un maillon essentiel du réseau universitaire de musique historique grâce auquel le chercheur peut retracer tout le processus de création d’une œuvre musicale, depuis la source patrimoniale conservée en bibliothèque jusqu’à la scène, grâce à des plateformes spécialisées à l'instar de LaDocumenta.eu. 


\subsection*{Étude de cas : le consortium européen LaDocumenta.eu}
\phantomsection\label{subsec:ladocumenta_cas} %pointe le chapitre pour permettre un renvoi d'un chapitre précédent vers ce chapitre.
Le consortium européen LaDocumenta.eu offre un terrain intéressant pour étudier les bibliothèques d'orchestre en Europe. Composé de cinq ensembles musicaux que sont \textit{Insula Orchestra} (France), le Collegium 1704 (République Tchèque), Le Concentus Musicus de Wien (Autriche), le Balthasar Neumann (Allemagne) et l'Accademia Bizantina (Italie), il est mû par la volonté de créer une plateforme numérique européenne pour la gestion des archives et la sauvegarde du patrimoine musical de la pratique historiquement informée (HIP). 
Ainsi, ce consortium montre qu'il existe un réseau international (ici européen) qui relie les bibliothèques musicales. La dimension numérique du consortium, reliant toutes les bibliothèques d'orchestre de la plateforme, renforce le travail en réseau et enterre cette légende désuète d'une bibliothèque vieillissante repoussant l'innovation. Au contraire, le réseau européen mobilisé pour ce projet regroupant l'IReMus et le Nikolaus Harnoncourt Zentrum permet de proposer une plateforme tournée vers la science ouverte. Adossés à des universités de référence (la Sorbonne pour l'IReMus et l'Université de Linz pour le Nikolaus Harnoncourt Zentrum), ces centres permettent de lier le projet à la recherche scientifique en musicologie afin de contextualiser les ressources de la plateforme. 
Pensée par les chefs pour répondre à des besoins concrets des orchestres, cette plateforme est également pleinement ancrée dans la science ouverte et le partage des ressources. L'objectif est également de conserver et de contextualiser les sources de la pratique historiquement informée (HIP) de la musique afin de développer la mémoire de cette pratique artistique et d'imaginer un jour sa reconnaissance comme patrimoine culturel immatériel de l'UNESCO.
